\chapter{Projective comparison with Boucksom--Jonsson theory}
\index{comparison!DXZ--BJ}\index{projective comparison theorem}
\label{chap:dxz-bj-comparison}

There are now two projective non-Archimedean psh spaces.
One is built from inverse systems of \(\cI\)-model test curves in K\"ahler
perturbations of a pseudo-effective class.
The other consists of extended-real psh metrics on the trivially valued
Berkovich analytification, defined by decreasing generically finite
Fubini--Study approximation.
The analytification constructed in the previous chapter identifies them.

The exact comparison theorem is due to Darvas--Xia--Zhang
\cite{DXZ23}.
Compatibility with operations is developed in
\cite{Xia23Operations}.
We keep track of which parts are formal and which use the
Boucksom--Jonsson envelope and approximation theory.

Throughout this chapter,
\[
\begin{gathered}
 X\text{ is a connected projective manifold over }\mathbf C,\\
 L\text{ is a pseudo-effective }\mathbf Q\text{-line bundle},\\
 h\text{ is a smooth metric on }L,\qquad \theta=c_1(L,h).
\end{gathered}
\tag{57.1}\label{eq:comp-scope}
\]
The core theorem below is not stated here for singular \(X\), a merely
compact K\"ahler \(X\), or an arbitrary transcendental class.

The comparison keeps four normalizations aligned.  A test-curve endpoint
is the value at the trivial valuation; parameter translation is addition
of a constant on \(X^{\An}\); an \(\mathcal I\)-envelope does not change
the valuative function; and the order on test curves becomes pointwise
order.  These checks prevent the bijection from being only an unnormalized
correspondence of rays.

\section{The comparison map}

Let \(\PSH^{\NA}_{\mathrm{tc}}(X,\theta)\) denote the test-curve space, and
let \(\PSH(L^{\An})\) denote the Boucksom--Jonsson space.
Here
\[
 \PSH^{\NA}_{\mathrm{tc}}(X,\theta)
 :=\PSH^{\NA}(X,\theta)
\]
is the DXZ inverse-limit space of
\Cref{def:tna-pshna}; the subscript only distinguishes this presentation
from the function-valued BJ presentation and introduces no third space.
\Cref{chap:valuative-analytification} constructed
\[
 \mathfrak A:
 \PSH^{\NA}_{\mathrm{tc}}(X,\theta)
       \longrightarrow\PSH(L^{\An}),
 \qquad
 \Gamma\longmapsto\Gamma^{\An}.
\tag{57.2}\label{eq:comp-map}
\]
On divisorial valuations,
\[
 \Gamma^{\An}(v)
 =\sup_{\tau\in\mathbf R}
       \left\{(\Gamma_\tau^{\theta+\omega})^{\An}(v)+\tau\right\},
\tag{57.3}\label{eq:comp-legendre}
\]
where \(\omega\) is any K\"ahler form.
The \(\cI\)-projection transition maps make this independent of
\(\omega\).
At the trivial valuation,
\[
                       \Gamma^{\An}(v_{\triv})=\Gamma_{\max}.
\tag{57.4}\label{eq:comp-max}
\]

\begin{proposition}\label{prop:comp-order-reflection}
For \(\Gamma,\Lambda\in\PSH^{\NA}_{\mathrm{tc}}(X,\theta)\),
\[
            \Gamma\leq\Lambda
       \quad\Longleftrightarrow\quad
            \Gamma^{\An}\leq\Lambda^{\An}
            \text{ on }X^{\An}.
\tag{57.5}\label{eq:comp-order-equivalence}
\]
In particular, \(\mathfrak A\) is injective.
\end{proposition}

\begin{proof}
The forward implication follows term by term from
\eqref{eq:comp-legendre}.

Conversely, suppose
\(\Gamma^{\An}\leq\Lambda^{\An}\).
Fix \(v\in X^{\Div}\) and \(t>0\).
Homogeneity gives
\[
 t\Gamma^{\An}(t^{-1}v)
   =\sup_{\tau\in\mathbf R}
       \left\{
       (\Gamma_\tau^{\theta+\omega})^{\An}(v)+t\tau
       \right\}.
\tag{57.6}\label{eq:comp-scaled-legendre}
\]
The assumed pointwise inequality therefore gives an inequality between the
two right sides for every \(t>0\).  For fixed \(v\), the functions
\(\tau\mapsto-(\Gamma_\tau^{\theta+\omega})^{\An}(v)\) and its
\(\Lambda\)-analogue are closed convex functions on their effective
intervals.  Legendre--Fenchel biconjugacy gives
\[
 (\Gamma_\tau^{\theta+\omega})^{\An}(v)
 \leq
 (\Lambda_\tau^{\theta+\omega})^{\An}(v)
\tag{57.7}\label{eq:comp-level-valuations}
\]
for every \(\tau\) and every divisorial \(v\).

The divisorial multiplier-ideal criterion turns this into
\[
 \Gamma_\tau^{\theta+\omega}
       \preceq_{\cI}
 \Lambda_\tau^{\theta+\omega}.
\]
Both levels are \(\cI\)-model potentials.
Applying the monotone \(\cI\)-envelope to representatives of the displayed
\(\cI\)-order shows that an \(\cI\)-model is the maximal normalized
representative of its constraint.  Hence
\[
 \Gamma_\tau^{\theta+\omega}
       \leq\Lambda_\tau^{\theta+\omega}.
\]
This is exactly \(\Gamma\leq\Lambda\).
\end{proof}

No envelope theorem is used in
\Cref{prop:comp-order-reflection}.
Only homogeneous Legendre inversion, the valuative multiplier-ideal
criterion, and maximality of \(\cI\)-model levels enter.

\section{The generically finite Fubini--Study core}

A rational generically finite Fubini--Study metric has the form
\[
 \phi=\frac1m\max_{1\leq j\leq M}
               \{\log|s_j|+\lambda_j\},
\tag{57.8}\label{eq:comp-gf}
\]
where \(mL\) is integral,
\(s_j\in H^0(X,mL)\setminus\{0\}\), and
\(\lambda_j\in\mathbf Q\).
The sections need not have empty common zero locus.
At a semivaluation supported in that locus, \(\phi=-\infty\).

\begin{lemma}\label{lem:comp-gf-realization}
Let \(T\) be an ample \(\mathbf Q\)-line bundle with a smooth positive
curvature form \(\omega_T\).  Every rational GF metric \(\phi\) of \(L\)
is \((\Gamma^T)^{\An}\) for a positive-mass
\(\mathcal I\)-model test curve \(\Gamma^T\) in the class
\(\theta+\omega_T\).  Its levels can be chosen elementary.
\end{lemma}

\begin{proof}
Group the indices in \eqref{eq:comp-gf} by the finitely many values
\(\tau_i=m^{-1}\lambda_j\).
At level \(\tau_i\), set
\[
 A_{\tau_i}
   :=\frac1m\max_{\lambda_j/m=\tau_i}
                    \log|s_j|_{h^m}^2.
\tag{57.9}\label{eq:comp-elementary-level}
\]
Interpolate affinely, make the curve constant below the smallest level, and
\(-\infty\) above the largest.
The PL formula gives
\[
 A^{\An}
  =\max_i\{A_{\tau_i}^{\An}+\tau_i\}
  =\phi.
\tag{57.10}\label{eq:comp-gf-pl}
\]
Upper concavification followed by levelwise
\(\mathcal I\)-projection in the class \(\theta+\omega_T\) produces a
genuine \(\Gamma^T\) without changing its divisorial analytification.
Divisorial determination upgrades equality to all of \(X^{\An}\).

It remains to check positive mass.  Before projection, the
\(-\infty\)-level \(A_{-\infty}\) is \(\theta\)-psh, hence
\[
 \theta+\omega_T+\ddc A_{-\infty}\geq\omega_T.
\]
The ordinary model projection preserves its mass, while the
\(\mathcal I\)-model projection is less singular and cannot decrease mass.
Consequently
\[
 \int_X\left(\theta+\omega_T+
       \ddc\Gamma^T_{-\infty}\right)^n
 \geq\int_X\omega_T^n>0.
\]
Thus \(\Gamma^T\) belongs to the required positive-mass test-curve space.

The factors \(m^{-1}\) occur both in the complex potential
\eqref{eq:comp-elementary-level} and in the BJ section function
\eqref{eq:comp-gf}.  Since our complex convention uses squared norms and
normalizes \(\nu_E(\log|s|_{h^m}^2)=\ord_E(s)\), analytification gives
\(m^{-1}\log|s|\) with no additional factor of two.

If all \(s_j\) vanish on the support of \(v\), every term in
\eqref{eq:comp-gf} and \eqref{eq:comp-gf-pl} is \(-\infty\).
Thus proper-support behavior is preserved.
\end{proof}

\begin{example}[Two weighted sections on \(\mathbf P^1\)]
Let \(L=\mathcal O_{\mathbf P^1}(1)\), write \([x:y]\) for homogeneous
coordinates, and fix \(a\in\mathbf Q_{>0}\).  The rational GF metric
\[
                 \phi=\max\{\log|x|,\log|y|+a\}
\]
has no common \(-\infty\)-locus because \(x\) and \(y\) do not vanish
simultaneously.  The construction in
\Cref{lem:comp-gf-realization} uses the two parameter levels \(a\) and
\(0\), interpolates between their single-section potentials, and then
takes the concave \(\cI\)-model replacement in
\(c_1(L+T)\), for any chosen ample \(T\).  Equation
\eqref{eq:comp-gf-pl} gives a positive-mass test curve \(\Gamma^T\) with
\[
                         (\Gamma^T)^{\An}=\phi.
\]

The identity can be checked on the two divisorial rays.  For
\(v_0=c\,\ord_{\{x=0\}}\) and
\(v_\infty=c\,\ord_{\{y=0\}}\), with \(c>0\), one obtains
\[
       \phi(v_0)=a,
       \qquad
       \phi(v_\infty)=\max\{0,a-c\}.
\]
At the trivial valuation, \(\phi(v_{\triv})=a\), so
\((\Gamma^T)_{\max}=a\) by \eqref{eq:comp-max}.  This elementary example checks
both the sign in the Legendre term and the normalization of the terminal
parameter.
\end{example}

\section{Decreasing limits}

The test-curve infimum of a decreasing net is defined under
\[
                         \inf_i\Gamma^i_{\max}>-\infty.
\tag{57.11}\label{eq:comp-noncollapse}
\]
This condition is visible on the BJ side by \eqref{eq:comp-max}.

\begin{proposition}\label{prop:comp-decreasing-nets}
Let \((\Gamma^i)_{i\in I}\) be a decreasing net satisfying
\eqref{eq:comp-noncollapse}.
Then
\[
                 \left(\inf_i\Gamma^i\right)^{\An}
                    =\inf_i(\Gamma^i)^{\An}
\tag{57.12}\label{eq:comp-decreasing}
\]
pointwise on \(X^{\An}\).
\end{proposition}

\begin{proof}
The componentwise infimum is legitimate because all levels below the
common terminal parameter have positive mass; the required envelope
identity is \cite[Proposition~2.8]{Xia23Operations}.  More directly,
\cite[Lemma~4.10]{Xia23Operations} states for a decreasing net that
\[
 \left(\inf_i\Gamma^i\right)^{\An}(v)
       =\inf_i(\Gamma^i)^{\An}(v)
\]
for every divisorial valuation \(v\), and also identifies the terminal
parameter.  The decreasing limit on the right is a BJ psh metric under
\eqref{eq:comp-noncollapse}.  Divisorial determination then upgrades the
identity from divisorial valuations to all of \(X^{\An}\).  This avoids
any unsourced minimax interchange for an arbitrary directed set.
\end{proof}

No equality
\(\vol(\inf_i\Gamma^i)=\lim_i\vol(\Gamma^i)\)
is assumed.
That extra hypothesis appeared for DH continuity because convergence of the
base convex bodies was needed there.

\section{The bijection}

\begin{theorem}\label{thm:comp-dxz}
Under \eqref{eq:comp-scope}, analytification is an order isomorphism
\[
 \mathfrak A:
 \PSH^{\NA}_{\mathrm{tc}}(X,\theta)
       \xrightarrow{\ \sim\ }\PSH(L^{\An}).
\tag{57.13}\label{eq:comp-bijection}
\]
\end{theorem}

\begin{proof}
Injectivity is \Cref{prop:comp-order-reflection}.
Let \(\phi\in\PSH(L^{\An})\).
The only nonformal surjectivity input is
\cite[Theorem~3.12(ii)]{DXZ23}: if \(T\) is ample, every
\(\phi\in\PSH(L^{\An})\), viewed as an element of
\(\PSH((L+T)^{\An})\), is the analytification of a positive-mass
\(\mathcal I\)-model test curve in \(c_1(L+T)\).  Its proof first treats
rational GF metrics by the construction of
\Cref{lem:comp-gf-realization}, then uses decreasing GF approximation with
a uniform positive-mass bound.  We import that realization theorem rather
than reproduce its analytic compactness argument.

Call a K\"ahler form rational if its class is the first Chern class of an
ample \(\mathbf Q\)-line bundle.  For every rational K\"ahler form
\(\omega\), the imported theorem gives a curve
\[
 \Gamma^\omega\in
 \PSH^{\NA}_{\mathrm{tc}}(X,\theta+\omega)_{>0},
 \qquad
                         (\Gamma^\omega)^{\An}=\phi.
\tag{57.14}\label{eq:comp-rational-component}
\]
If \(\omega\leq\omega'\) are rational, projection invariance gives
\[
 \left(P_{\theta+\omega'}[\Gamma^\omega]_{\mathcal I}\right)^{\An}
 =\phi=(\Gamma^{\omega'})^{\An}.
\]
Injectivity in the big class \(\theta+\omega'\) therefore yields
\[
 P_{\theta+\omega'}[\Gamma^\omega]_{\mathcal I}
 =\Gamma^{\omega'}.
\tag{57.15}\label{eq:comp-rational-cocycle}
\]

Now let \(\eta\) be an arbitrary K\"ahler form.  Choose a rational
K\"ahler form \(\omega\leq\eta\) and set
\[
 \Gamma^\eta
 :=P_{\theta+\eta}[\Gamma^\omega]_{\mathcal I}.
\tag{57.16}\label{eq:comp-arbitrary-component}
\]
This is independent of \(\omega\).  Indeed, for two choices
\(\omega_1,\omega_2\leq\eta\), choose a sufficiently small positive
rational multiple \(\omega_0\) of one fixed rational K\"ahler form so that
\(\omega_0\leq\omega_1,\omega_2\).  Applying
\eqref{eq:comp-rational-cocycle} twice shows that both candidates in
\eqref{eq:comp-arbitrary-component} are
\(P_{\theta+\eta}[\Gamma^{\omega_0}]_{\mathcal I}\).
The same argument with a common rational lower form proves the transition
cocycle for arbitrary \(\eta\leq\eta'\).
Thus \((\Gamma^\eta)_\eta\) is a compatible inverse system.  Projection
invariance gives \((\Gamma^\eta)^{\An}=\phi\) for every component, hence
\(\Gamma^{\An}=\phi\).  This is the surjectivity argument of
\cite[Theorem~3.14]{DXZ23}.
\end{proof}

\begin{remark}[Where envelopes enter]
\label{rem:comp-envelope-dependence}
The map, its well-definedness, order reflection, GF realization, and
decreasing-limit identity were proved without surjectivity.
Surjectivity uses the realization theorem
\cite[Theorem~3.12(ii)]{DXZ23}, whose proof uses sequential GF
regularization from the Boucksom--Jonsson approximation and envelope
theory.
Thus bijectivity is not a formal consequence of the multiplier-ideal
formula.

For the smooth projective \(X\) considered here, the needed input is
available in the cited BJ theory.
The general unibranch envelope theorem proved in
\Cref{chap:bj-envelope-consequences} supplies the broader algebraic
envelope statement.
The hypotheses of a comparison construction must still be checked before
transporting the bijection to another category.
\end{remark}

\section{Operations and envelopes}

The order isomorphism transports order-theoretic operations.
The direct test-curve proofs in \cite{Xia23Operations} show that this
transport agrees with the operations already constructed on the
transcendental side.

\begin{theorem}\label{thm:comp-operations}
Whenever the displayed operations are defined,
\[
\begin{aligned}
 (\Gamma+C)^{\An}
     &=\Gamma^{\An}+C,\\
 (\lambda\Gamma)^{\An}
     &=\lambda\Gamma^{\An}\qquad(\lambda>0),\\
 (\Gamma+\Lambda)^{\An}
     &=\Gamma^{\An}+\Lambda^{\An},\\
 (\Gamma\vee\Lambda)^{\An}
     &=\max\{\Gamma^{\An},\Lambda^{\An}\}.
\end{aligned}
\tag{57.18}\label{eq:comp-elementary-operations}
\]
If \((\Gamma^i)\) is nonempty and
\(\sup_i\Gamma^i_{\max}<+\infty\), then
\[
 \left(\uscsup_i\Gamma^i\right)^{\An}
       =\left(\sup_i(\Gamma^i)^{\An}\right)^*.
\tag{57.19}\label{eq:comp-sup}
\]
For a decreasing net satisfying \eqref{eq:comp-noncollapse},
\eqref{eq:comp-decreasing} holds.
\end{theorem}

\begin{proof}
Translations and finite maxima follow directly from the Legendre formula.
Positive rescaling follows from
\eqref{eq:comp-scaled-legendre}.
Addition follows levelwise from addition of psh singularities and the
infimal-convolution formula for test curves; analytification of a sum of
complex singularities adds their divisorial values.

The divisorial maximum identity is
\cite[Corollary~4.8]{Xia23Operations}; addition and rescaling are
\cite[Lemma~4.17]{Xia23Operations}.  Since both sides of each identity are
BJ psh metrics, divisorial determination extends the equality to
\(X^{\An}\).  This is where the proved projective comparison is used; the
same sentence would not justify equality on arbitrary nondivisorial points
of the tropical spectrum of a nonprojective complex space.

For the supremum, the test-curve construction is the levelwise regularized
supremum followed by \(\cI\)-projection.
The projection preserves analytification, and the projective envelope
property makes the upper semicontinuous regularization on the right of
\eqref{eq:comp-sup} a BJ psh metric.
Order reflection identifies it with the left side.
The decreasing statement is
\Cref{prop:comp-decreasing-nets}.
\end{proof}

\begin{remark}[Envelope firewall]
\label{rem:comp-operation-envelope-firewall}
Equation \eqref{eq:comp-sup} is the operation in which the envelope theorem
is visible.
Without the appropriate envelope property, the regularized supremum of a
bounded-above psh family need not already be known to be psh.
Finite maxima, translations, sums, positive rescaling, and the admissible
decreasing infimum do not require that additional assertion.
\end{remark}

\section{Monotone limits and order topology}

The comparison identifies the natural monotone convergences, but the two
directions have different endpoint hypotheses.

\begin{proposition}\label{prop:comp-increasing-nets}
Let \((\Gamma^i)_{i\in I}\) be increasing and assume
\[
                         \sup_i\Gamma^i_{\max}<+\infty.
\]
Then
\[
 \Gamma:=\uscsup_i\Gamma^i
\]
is defined and
\[
 \Gamma^{\An}
       =\left(\sup_i(\Gamma^i)^{\An}\right)^*.
\tag{57.20}\label{eq:comp-increasing}
\]
Moreover,
\[
                         \vol(\Gamma)=\lim_i\vol(\Gamma^i).
\tag{57.21}\label{eq:comp-increasing-volume}
\]
\end{proposition}

\begin{proof}
The first identity is \eqref{eq:comp-sup}.
For the volume statement, project every metric to one fixed K\"ahler
perturbation.
The positive-mass test-curve theorem gives monotone continuity there.
The error introduced by a sufficiently small perturbation is uniformly
bounded by mixed masses with a fixed K\"ahler class.
Letting the perturbation tend to zero proves
\eqref{eq:comp-increasing-volume}.
The supremum identity itself is
\cite[Theorem~4.13]{Xia23Operations}; the uniform mixed-mass estimate is
the same one used in \Cref{thm:operations-volume}.
\end{proof}

\begin{remark}
For decreasing nets, the condition is the lower bound
\eqref{eq:comp-noncollapse}, and the volume may drop:
\[
                  \vol(\inf_i\Gamma^i)\leq\lim_i\vol(\Gamma^i).
\]
The reverse inequality fails in general.
Thus order convergence and energy or DH convergence must not be conflated.
\end{remark}

\section{Projection and envelopes}

For every K\"ahler form \(\omega\), the transition map satisfies
\[
 \left(P_{\theta+\omega}[\Gamma]_{\cI}\right)^{\An}
                   =\Gamma^{\An}.
\tag{57.22}\label{eq:comp-projection-invariance}
\]
This is not injectivity failure:
the source and target components live in different cohomology classes inside
the inverse system.
The comparison forgets the auxiliary perturbation while retaining the
common valuative metric.

\begin{proposition}\label{prop:comp-perturbed-characterization}
A compatible family
\[
 \Gamma^{\theta+\omega}\in
       \PSH^{\NA}_{\mathrm{tc}}(X,\theta+\omega)_{>0}
\]
defines a metric in class \(\theta\) if and only if for
\(\omega'\leq\omega''\),
\[
 P_{\theta+\omega''}
       [\Gamma^{\theta+\omega'}]_{\cI}
       =\Gamma^{\theta+\omega''}.
\tag{57.23}\label{eq:comp-cocycle}
\]
Under comparison, all components in such a family have one common
analytification.
\end{proposition}

\begin{proof}
The first assertion is the inverse-system definition.
Projection invariance gives equality of analytifications in the forward
direction.
Conversely, if the components have a common analytification, the two sides
of \eqref{eq:comp-cocycle} have the same analytification.
Order reflection, hence injectivity in the perturbed class, makes them equal.
\end{proof}

This proposition is the step used in the surjectivity proof after constructing
\eqref{eq:comp-arbitrary-component}.
It is important that injectivity had already been proved without assuming
surjectivity.

\section{Trace and restriction}

Let \(Y\subseteq X\) be an irreducible analytic subset and
\(\widetilde Y\to Y\) its normalization.
Denote by \(v_{Y,\triv}\in X^{\An}\) the trivial valuation with support
\(Y\).

\begin{proposition}\label{prop:comp-restriction-criterion}
For \(\Gamma\in\PSH^{\NA}_{\mathrm{tc}}(X,\theta)\), the following are
equivalent:
\begin{enumerate}
\item \(\Gamma^{\An}(v_{Y,\triv})>-\infty\);
\item \(\Gamma^{\An}|_{Y^{\An}}\not\equiv-\infty\);
\item the test-curve trace \(\Tr_Y(\Gamma)\) is well-defined.
\end{enumerate}
\par
\end{proposition}

\begin{proof}
The maximum principle for BJ psh metrics
\cite[Lemma~1.4(i)]{BJ22} gives the equivalence of the first two conditions.
For a complex potential \(\psi\),
\[
 \psi^{\An}(v_{Y,\triv})
 =
 \begin{cases}
 -\infty,&\nu(\psi,Y)>0,\\
 0,&\nu(\psi,Y)=0.
 \end{cases}
\tag{57.24}\label{eq:comp-generic-trivial-value}
\]
Applying this to the levels in \eqref{eq:comp-legendre} shows that the first
condition holds exactly when sufficiently negative levels have zero generic
Lelong number along \(Y\), which is the definition of the trace domain.
\end{proof}

\begin{theorem}\label{thm:comp-trace}
Assume the equivalent conditions of
\Cref{prop:comp-restriction-criterion}.
Then
\[
 \Tr_Y(\Gamma)^{\An}|_{Y^{\Div}}
       =\Gamma^{\An}|_{Y^{\Div}}.
\tag{57.25}\label{eq:comp-trace}
\]
\end{theorem}

\begin{proof}
First apply an \(\cI\)-projection to a K\"ahler perturbation.
It changes neither side on divisorial valuations and reduces to positive
mass.

Choose a decreasing GF approximation
\(\phi_j\searrow\Gamma^{\An}\), realize each \(\phi_j\) by a PL test curve
\(\Gamma^j\), and project all curves to one fixed larger K\"ahler class.
\Cref{prop:comp-decreasing-nets} and injectivity show that the projected
curves decrease to the projection of \(\Gamma\).
Levelwise their \(\cI\)-model potentials converge in \(d_S\).
Continuity of the complex trace under decreasing \(d_S\)-limits gives, for
each divisorial \(v\) over \(Y\),
\[
 \Tr_Y(\Gamma)^{\An}(v)
       =\inf_j\Tr_Y(\Gamma^j)^{\An}(v).
\tag{57.26}\label{eq:comp-trace-limit}
\]

It remains to treat a GF metric.
Its realizing PL curve has levels that are finite maxima of logarithms of
sections.
When restriction is nontrivial, restricting these sections to
\(\widetilde Y\) commutes with divisorial order of vanishing.
Thus for every level \(A_\tau\),
\[
       \Tr_Y(A_\tau)^{\An}(v)=A_\tau^{\An}(v).
\]
Taking the finite Legendre maximum proves
\eqref{eq:comp-trace} for \(\Gamma^j\).
Passing to the infimum in \eqref{eq:comp-trace-limit} proves the theorem.
This gives exactly the divisorial trace formula of
\cite[Section~4.8]{Xia23Operations}; no nondivisorial restriction theorem
has been inserted in the limiting argument.
\end{proof}

The conclusion is deliberately stated on \(Y^{\Div}\), exactly as proved by
the trace comparison.
It is not a blanket assertion that restriction of every extended-real BJ
psh function is automatically a BJ psh metric on \(Y^{\An}\).
The nontriviality criterion and analytic trace are essential.

\section{Scope ledger}

The comparison has four separate layers:
\begin{enumerate}
\item the multiplier-ideal formula constructs \(\mathfrak A\) on every
mass stratum;
\item homogeneous Legendre inversion makes it an order embedding;
\item sequential GF approximation, an envelope-theoretic input, makes it
surjective in the projective rational setting;
\item the operation theorem and trace continuity identify the extra
structures.
\end{enumerate}

The cited paper \cite{DXZ23} also develops the transcendental inverse-system
theory beyond rational Hodge classes.
The exact BJ bijection \eqref{eq:comp-bijection}, however, uses the
projective algebraic analytification and is asserted here only under
\eqref{eq:comp-scope}.
No comparison with the tropical-spectrum psh theory of a general complex
space has been used.
That requires additional results reserved for later chapters.

For orientation, \Cref{thm:ts-algebraic-recovery} will identify the
underlying space \(X^{\An}\) with the tropical spectrum of coherent ideals
when \(X\) is proper.  Under that homeomorphism,
\(\log|\mathfrak a|=-v(\mathfrak a)\), so the multiplier-ideal formula of
\Cref{def:va-psh-analytification} is literally an ideal-function formula on
the tropical spectrum.  This underlying-space dictionary is formal from
GAGA and does not by itself compare positivity.  The additional positivity
argument in \Cref{thm:pseff-projective-full-mp-bj} proves that, when \(X\)
is connected and smooth projective and
\(\alpha\in\operatorname{Amp}(X)\subset N^1(X)_{\mathbf R}\), this
homeomorphism identifies the full Mesquita--Piccione and
Boucksom--Jonsson psh cones, without a finite-energy hypothesis.  If
\(\alpha=c_1(L)\) for an ample \(\mathbf Q\)-line bundle, composition with
\eqref{eq:comp-bijection} gives the full DXZ--MP order isomorphism on all
semivaluations.  Neither conclusion is asserted here for a nonprojective
compact K\"ahler manifold, a transcendental class outside
\(N^1(X)_{\mathbf R}\), or a non-nef class; equality of the psh cones also
does not by itself identify their Monge--Amp\`ere operators.

\begin{remark}[Primary-source boundary]
The order isomorphism is \cite[Theorem~3.14]{DXZ23}.  Its source-side
objects are DXZ inverse systems of \(\mathcal I\)-model test curves, and
its target is the BJ cone in the projective rational class.  The operation
formulas, endpoint-bounded monotone limits, and trace comparison are from
\cite{Xia23Operations}, with the exact theorem numbers recorded at their
uses above.  These source statements do not supply a compact-K\"ahler BJ
analytification or a singular-space comparison.  The full MP--BJ
identification in the smooth projective ample N\'eron--Severi setting is the
separate result \Cref{thm:pseff-projective-full-mp-bj}; its hypotheses do
not enlarge the DXZ theorem proved in this chapter to the compact-K\"ahler,
transcendental, or non-nef settings.
\end{remark}
